\documentclass[12pt]{article}
\usepackage{fullpage}
\usepackage[T1]{fontenc}
\usepackage[utf8]{inputenc}
\usepackage{lmodern}
\usepackage{microtype}
\usepackage{amsmath,amssymb}
\usepackage{amsthm}
\usepackage{graphicx}
\usepackage{booktabs}
\usepackage{hyperref}
\usepackage{url}
\usepackage{color}
\usepackage{xcolor}
\usepackage{enumitem}
\usepackage{amsfonts}
\usepackage{algorithm}
\usepackage{algpseudocode}

\DeclareMathOperator{\vecop}{vec}
\DeclareMathOperator{\diag}{diag}
\newcommand{\bR}{\mathbb{R}}
\newcommand{\KR}{\odot}

\newtheorem{theorem}{Theorem}
\newtheorem{lemma}{Lemma}
\newtheorem{proposition}{Proposition}

\title{Problem 10: Preconditioned CG for RKHS-Constrained CP Decomposition}
\author{}
\date{}

\begin{document}
\maketitle

\section{Setup and notation}

We recall the problem data. Mode $k$ has size $n\equiv n_k$, the remaining modes
multiply to $M=\prod_{i\neq k}n_i$, the full tensor has $N=\prod_i n_i = nM$
entries, of which $q\ll N$ are observed. We assume throughout
\[
  d<n,\ r<q\ll N,\qquad n\ll M ,
\]
where $d$ is the order of the tensor and $r$ is the CP rank.
The matrix $T\in\bR^{n\times M}$ is the mode-$k$ unfolding (missing entries
set to $0$), $Z=A_d\KR\cdots\KR A_{k+1}\KR A_{k-1}\KR\cdots\KR A_1\in\bR^{M\times r}$
is the Khatri--Rao product of the fixed factors, $K\in\bR^{n\times n}$ is the
symmetric positive semidefinite RKHS kernel matrix for mode~$k$, and
$B=TZ\in\bR^{n\times r}$ is the MTTKRP. The selection matrix
$S\in\bR^{N\times q}$ is a column subset of $I_N$ such that $S^{\mathsf T}\vecop(T)$
extracts the $q$ observed entries from $\vecop(T)$. We must solve, for the
unknown $W\in\bR^{n\times r}$,
\begin{equation}\label{eq:system}
  G\,\vecop(W) = (I_r\otimes K)\,\vecop(B),
  \qquad
  G := (Z\otimes K)^{\mathsf T}SS^{\mathsf T}(Z\otimes K)+\lambda\,(I_r\otimes K),
\end{equation}
a symmetric system of order $nr$. (The CP factor of mode $k$ is then $A_k=KW$.)
We use $\vecop$ in the standard column-stacking convention, and we order the
unknown so that the entry of $\vecop(W)$ in position $p\cdot n+s$ (with
$0\le p<r$, $0\le s<n$) equals $W[s,p]$; with this convention the Kronecker
products in \eqref{eq:system} use the ordering
$(Z\otimes K)[\,m\cdot n+i,\ p\cdot n+s\,]=Z[m,p]\,K[i,s]$, consistent with
$\vecop(T)$ being the column stacking of the $n\times M$ matrix $T$. All claims
below are invariant under the choice of consistent ordering.

\paragraph{Encoding the observations.}
Each observed entry is a pair $(i_\ell,m_\ell)\in\{0,\dots,n-1\}\times\{0,\dots,M-1\}$,
$\ell=1,\dots,q$, namely the row and column of $T$ that the $\ell$-th column of
$S$ selects. Define the \emph{subsampled} factor blocks
\begin{equation}\label{eq:subsample}
  \tilde Z\in\bR^{q\times r},\quad \tilde Z[\ell,:]=Z[m_\ell,:],
  \qquad
  \tilde K\in\bR^{q\times n},\quad \tilde K[\ell,:]=K[i_\ell,:].
\end{equation}
These are exactly the ``subset of rows of $Z$ and $K$'' referred to in the
statement, gathered into $q\times r$ and $q\times n$ matrices. Forming
\eqref{eq:subsample} costs $O(q\,r)$ and $O(q\,n)$ time and $O(q(r+n))$ storage;
no object of order $M$ or $N$ is created.

The whole development rests on the following exact factorization of the data
block of $G$.

\begin{lemma}[Subsampled Khatri--Rao factorization]\label{lem:P}
Let $P:=S^{\mathsf T}(Z\otimes K)\in\bR^{q\times nr}$. Then row $\ell$ of $P$ is
$Z[m_\ell,:]\otimes K[i_\ell,:]=\tilde Z[\ell,:]\otimes\tilde K[\ell,:]$, i.e.
$P=\tilde Z\KR_{\mathrm{row}}\tilde K$ is the row-wise (transposed) Khatri--Rao
product of $\tilde Z$ and $\tilde K$. Consequently
\begin{equation}\label{eq:G-PtP}
  G=P^{\mathsf T}P+\lambda\,(I_r\otimes K).
\end{equation}
\end{lemma}

\begin{proof}
$S^{\mathsf T}$ selects $q$ rows of $Z\otimes K$, namely the rows indexed by the
observed flat indices $i_\ell+n\,m_\ell$. By the Kronecker ordering above,
row $m_\ell\cdot n+i_\ell$ of $Z\otimes K$ equals the Kronecker product of
row $m_\ell$ of $Z$ with row $i_\ell$ of $K$, which by \eqref{eq:subsample} is
$\tilde Z[\ell,:]\otimes\tilde K[\ell,:]$. Stacking the rows gives
$P=\tilde Z\KR_{\mathrm{row}}\tilde K$. Finally, since
$(SS^{\mathsf T})^{\mathsf T}=SS^{\mathsf T}$ and $S^{\mathsf T}=S^{\mathsf T}$,
$$(Z\otimes K)^{\mathsf T}SS^{\mathsf T}(Z\otimes K)
=\big(S^{\mathsf T}(Z\otimes K)\big)^{\mathsf T}\big(S^{\mathsf T}(Z\otimes K)\big)
=P^{\mathsf T}P,$$
which gives \eqref{eq:G-PtP}.
\end{proof}

Because $S^{\mathsf T}S=I_q$ ($S$ is a column subset of the identity) the middle
factor $SS^{\mathsf T}$ is an orthogonal projector, and Lemma~\ref{lem:P} shows
$G$ never has to touch a vector of length $M$ or $N$: all that survives the
projection are the $q$ subsampled rows.

\section{Properties of $G$}

\begin{proposition}\label{prop:spd}
$G$ is symmetric positive semidefinite. If $K\succ0$ and $\lambda>0$ then
$G\succ0$, and \eqref{eq:system} has a unique solution.
\end{proposition}

\begin{proof}
By \eqref{eq:G-PtP}, $G=P^{\mathsf T}P+\lambda(I_r\otimes K)$. The first summand
is a Gram matrix, hence symmetric PSD. For the second, $I_r\otimes K$ is
symmetric and its eigenvalues are products of eigenvalues of $I_r$ (all $1$) and
of $K$ (all $\ge0$); hence $I_r\otimes K\succeq0$, and $\succ0$ iff $K\succ0$.
A sum of a PSD and a positive (semi)definite matrix is positive (semi)definite,
giving the claims. Uniqueness follows since $G\succ0$ is nonsingular.
\end{proof}

In the RKHS setting $K$ is a kernel Gram matrix on the (few) mode-$k$ grid
points and is taken positive definite (a positive-definite kernel, or a PSD
kernel with a small nugget $K\leftarrow K+\varepsilon I_n$, $\varepsilon>0$);
we assume $K\succ0$ and $\lambda>0$, so $G\succ0$ and conjugate gradients
applies.

\section{Matrix-free matrix--vector product}

CG only needs the action $W\mapsto \mathcal{G}(W)$, where $\mathcal{G}(W)$ is
the $n\times r$ matrix with $\vecop(\mathcal{G}(W))=G\,\vecop(W)$. We compute it
without forming $G$, $Z\otimes K$, $S$, or any vector of length $M$ or $N$.

\begin{lemma}[Matrix-valued matvec]\label{lem:matvec}
For $W\in\bR^{n\times r}$ set $Y=KW\in\bR^{n\times r}$. Define
$u\in\bR^q$ by
\begin{equation}\label{eq:Pw}
  u_\ell=\sum_{p=0}^{r-1}\tilde Z[\ell,p]\,Y[i_\ell,p]
  \qquad(\ell=1,\dots,q),
\end{equation}
i.e. $u_\ell$ is the inner product of row $\ell$ of $\tilde Z$ with row $i_\ell$
of $Y$. Define $C\in\bR^{n\times r}$ by scatter-accumulation
\begin{equation}\label{eq:Ptu}
  C[i,:]=\sum_{\ell:\,i_\ell=i} u_\ell\,\tilde Z[\ell,:]
  \qquad(i=0,\dots,n-1).
\end{equation}
Then
\begin{equation}\label{eq:matvec}
  \mathcal{G}(W)=K\,C+\lambda\,K\,W=K\,(C+\lambda W).
\end{equation}
\end{lemma}

\begin{proof}
By Lemma~\ref{lem:P}, $P\,\vecop(W)$ has $\ell$-th entry
$(\tilde Z[\ell,:]\otimes\tilde K[\ell,:])\,\vecop(W)
=\sum_{p,s}\tilde Z[\ell,p]\,\tilde K[\ell,s]\,W[s,p]
=\sum_p\tilde Z[\ell,p]\sum_s K[i_\ell,s]W[s,p]
=\sum_p\tilde Z[\ell,p]\,(KW)[i_\ell,p]
=\sum_p\tilde Z[\ell,p]\,Y[i_\ell,p]=u_\ell$,
which is \eqref{eq:Pw}. Next, for the transpose,
$P^{\mathsf T}u$ has entry in position $p\cdot n+s$ equal to
$\sum_\ell u_\ell\,\tilde Z[\ell,p]\,K[i_\ell,s]
=\sum_s' \big(\sum_\ell u_\ell\,\tilde Z[\ell,p]\,[i_\ell=i]\big)K[i,s]$.
Writing $C[i,p]=\sum_{\ell:i_\ell=i}u_\ell\tilde Z[\ell,p]$ (which is
\eqref{eq:Ptu}) and using symmetry of $K$, the $(s,p)$ entry of the matrix form
of $P^{\mathsf T}u$ is $\sum_i K[s,i]C[i,p]=(KC)[s,p]$, i.e. the matrix form of
$P^{\mathsf T}P\,\vecop(W)$ is $KC$. Finally
$(I_r\otimes K)\vecop(W)=\vecop(KW)$, so the regularizer contributes
$\lambda KW$. Adding the two and using \eqref{eq:G-PtP} gives
\eqref{eq:matvec}.
\end{proof}

\paragraph{Cost of one matvec.}
Step $Y=KW$: $O(n^2r)$. Step \eqref{eq:Pw}: each $u_\ell$ is an $r$-dimensional
dot product, total $O(qr)$. Step \eqref{eq:Ptu}: each observation adds a scaled
length-$r$ row, total $O(qr)$. Step \eqref{eq:matvec}: one more $K$-multiply,
$O(n^2r)$. Hence one matvec costs
\begin{equation}\label{eq:matveccost}
  O(n^2 r + q r)
\end{equation}
time and $O(nr+q(n+r))$ memory, with \emph{no} dependence on $M$ or $N$.

\section{The preconditioner}

The dominant ill-conditioning of $G$ comes from the kernel factor: smooth RKHS
kernels have rapidly (often exponentially) decaying spectra, so $K$ is severely
ill-conditioned, and \emph{both} terms of $G=P^{\mathsf T}P+\lambda(I_r\otimes K)$
inherit that conditioning (the regularizer carries one factor of $K$, and each
rank-one term of $P^{\mathsf T}P$ carries the factor $\tilde K[\ell,:]\otimes\dots$
built from rows of $K$). We therefore precondition out the kernel factor with
\begin{equation}\label{eq:precond}
  \mathcal M := I_r\otimes K\ \succ0 .
\end{equation}
This is symmetric positive definite (Proposition~\ref{prop:spd} argument with
$\lambda=1,P=0$), so split-preconditioned CG (equivalently, CG in the
$\mathcal M$-inner product) is well defined.

\paragraph{Why this is the right choice.}
With $\mathcal M=I_r\otimes K$, the preconditioned operator is
\[
  \mathcal M^{-1}G
  =(I_r\otimes K^{-1})P^{\mathsf T}P+\lambda I_{nr}.
\]
Its spectrum is $\lambda+\sigma$ with $\sigma\ge0$ the eigenvalues of the PSD
matrix $\mathcal M^{-1/2}P^{\mathsf T}P\,\mathcal M^{-1/2}
=(I_r\otimes K^{-1/2})P^{\mathsf T}P(I_r\otimes K^{-1/2})$. The regularizer is
mapped exactly to the identity (so its enormous condition number is removed
completely), and the data term contributes a matrix of rank at most $q$ whose
nonzero eigenvalues are bounded by
$\|(I_r\otimes K^{-1/2})P^{\mathsf T}\|_2^2$. Thus
$\mathrm{cond}(\mathcal M^{-1}G)\le 1+\tfrac1\lambda\lambda_{\max}
\big(\mathcal M^{-1/2}P^{\mathsf T}P\mathcal M^{-1/2}\big)$, which is governed by
the data fit, not by the kernel spectrum. Empirically (see the verification note
below) this collapses the condition number from $\sim\!10^{15}$ to $\sim\!10^{4}$
on a smooth-kernel instance, and CG iteration counts drop several-fold.

\paragraph{Why applying $\mathcal M^{-1}$ is cheap.}
$\mathcal M^{-1}=I_r\otimes K^{-1}$, so solving $\mathcal M\,V=R$ for the
$n\times r$ matrix $R$ amounts to solving $K\,V=R$ column by column. Precompute a
Cholesky factorization $K=LL^{\mathsf T}$ \emph{once} in $O(n^3)$; then each
preconditioner application is $r$ pairs of triangular solves with $L,L^{\mathsf T}$,
costing $O(n^2 r)$ and never touching $M$ or $N$. (If $K$ is only PSD, factor
$K+\varepsilon I_n$ for a small nugget $\varepsilon>0$; this is the same matrix
used to enforce $G\succ0$.)

\paragraph{Right-hand side.}
The right-hand side is $(I_r\otimes K)\vecop(B)=\vecop(KB)$. The MTTKRP
$B=TZ\in\bR^{n\times r}$ is computed from the observed entries only: writing
$t_\ell$ for the value of the $\ell$-th observed entry $(i_\ell,m_\ell)$,
$B[i,:]=\sum_{\ell:i_\ell=i}t_\ell\,\tilde Z[\ell,:]$, a scatter-accumulation
costing $O(qr)$ (no order-$M$ or order-$N$ work). Then $KB$ costs $O(n^2r)$.

\section{Algorithms}

\begin{algorithm}[h]
\caption{Direct symmetric solve of the mode-$k$ subproblem \eqref{eq:system}}
\label{alg:direct}
\begin{algorithmic}[1]
\Require Subsampled rows $\tilde Z\in\bR^{q\times r}$, $\tilde K\in\bR^{q\times n}$
  and indices $\{(i_\ell,m_\ell)\}_{\ell=1}^q$; kernel $K\in\bR^{n\times n}$ (SPD);
  observed values $\{t_\ell\}$; regularizer $\lambda>0$.
\Ensure $W\in\bR^{n\times r}$ solving \eqref{eq:system}; factor $A_k=KW$.
\State Form $P\in\bR^{q\times nr}$ with $P[\ell,:]=\tilde Z[\ell,:]\otimes\tilde K[\ell,:]$
  \Comment{$O(q\,nr)$}
\State $G\gets P^{\mathsf T}P+\lambda\,(I_r\otimes K)$
  \Comment{Gram: $O(q\,n^2r^2)$; add regularizer: $O(n^2r)$}
\State $B[i,:]\gets\sum_{\ell:\,i_\ell=i} t_\ell\,\tilde Z[\ell,:]$ for all $i$;\quad $b\gets\vecop(KB)$
  \Comment{$O(qr+n^2r)$}
\State Compute Cholesky $G=R^{\mathsf T}R$ \Comment{$O(n^3r^3)$}
\State Solve $R^{\mathsf T}y=b$, then $R\,w=y$; set $W\gets$ matrix form of $w$
  \Comment{$O(n^2r^2)$}
\State $A_k\gets KW$ \Comment{$O(n^2r)$}
\State \Return $W,\ A_k$
\end{algorithmic}
\end{algorithm}

The direct method is dominated by forming and factoring the dense $nr\times nr$
matrix $G$: assembly $O(q\,n^2r^2)$ and Cholesky $O(n^3r^3)$, with $O(n^2r^2)$
memory. This is the ``standard linear solver'' cost the statement seeks to beat.

\begin{algorithm}[h]
\caption{Matrix-free Preconditioned CG for \eqref{eq:system}}
\label{alg:pcg}
\begin{algorithmic}[1]
\Require Subsampled rows $\tilde Z\in\bR^{q\times r}$ and indices
  $\{(i_\ell,m_\ell)\}_{\ell=1}^q$; kernel $K\in\bR^{n\times n}$ (SPD); observed
  values $\{t_\ell\}$; $\lambda>0$; tolerance $\tau$; max iterations $T_{\max}$.
\Ensure $W\in\bR^{n\times r}$ approximately solving \eqref{eq:system}; $A_k=KW$.
\Statex \textbf{Precomputation.}
\State Cholesky $K=LL^{\mathsf T}$ \Comment{$O(n^3)$, once}
\State $B[i,:]\gets\sum_{\ell:\,i_\ell=i}t_\ell\,\tilde Z[\ell,:]$;\quad $\mathrm{rhs}\gets KB$
  \Comment{$O(qr+n^2r)$}
\Statex \textbf{Operators (used as subroutines, never assembled).}
\Function{Matvec}{$W$}\Comment{returns $\mathcal G(W)$, Lemma~\ref{lem:matvec}}
  \State $Y\gets KW$;\quad $u_\ell\gets\sum_p\tilde Z[\ell,p]\,Y[i_\ell,p]\ \forall\ell$
  \State $C\gets 0$; \textbf{for} $\ell=1$ \textbf{to} $q$ \textbf{do} $C[i_\ell,:]\mathrel{+}=u_\ell\,\tilde Z[\ell,:]$
  \State \Return $K\,(C+\lambda W)$
\EndFunction
\Function{Prec}{$R$}\Comment{returns $\mathcal M^{-1}R=K^{-1}R$ columnwise}
  \State Solve $L\,\Xi=R$; solve $L^{\mathsf T}V=\Xi$; \Return $V$
\EndFunction
\Statex \textbf{PCG iteration.}
\State $W\gets 0$;\quad $R\gets\mathrm{rhs}-\Call{Matvec}{W}$
\State $V\gets\Call{Prec}{R}$;\quad $D\gets V$;\quad $\rho\gets\langle R,V\rangle$;\quad $\rho_0\gets\langle R,R\rangle$
\For{$t=1$ \textbf{to} $T_{\max}$}
  \State $Q\gets\Call{Matvec}{D}$;\quad $\alpha\gets\rho/\langle D,Q\rangle$
  \State $W\gets W+\alpha D$;\quad $R\gets R-\alpha Q$
  \If{$\langle R,R\rangle\le\tau^2\rho_0$} \textbf{break}\EndIf
  \State $V\gets\Call{Prec}{R}$;\quad $\rho'\gets\langle R,V\rangle$;\quad $\beta\gets\rho'/\rho$
  \State $D\gets V+\beta D$;\quad $\rho\gets\rho'$
\EndFor
\State $A_k\gets KW$ \Comment{$O(n^2r)$}
\State \Return $W,\ A_k$
\end{algorithmic}
\end{algorithm}

Here $\langle X,Y\rangle=\sum_{s,p}X[s,p]Y[s,p]=\vecop(X)^{\mathsf T}\vecop(Y)$,
computed in $O(nr)$.

\section{Correctness}

\begin{theorem}\label{thm:correct}
Assume $K\succ0$ and $\lambda>0$. Then:
\begin{enumerate}[label=(\alph*)]
\item Algorithm~\ref{alg:direct} returns the exact solution $W$ of
\eqref{eq:system}.
\item In exact arithmetic, Algorithm~\ref{alg:pcg} is the conjugate gradient
method applied to \eqref{eq:system} with preconditioner $\mathcal M=I_r\otimes K$;
it terminates with the exact solution in at most $nr$ iterations and, with the
stopping test, returns $W$ with relative residual $\|G\vecop(W)-b\|_2\le
\tau\,\|b\|_2$ when it breaks. Both algorithms return the same $W$ (up to the
PCG tolerance), hence the same $A_k=KW$.
\end{enumerate}
\end{theorem}

\begin{proof}
\emph{(a)} By Lemma~\ref{lem:P}, the matrix $G$ assembled in
Algorithm~\ref{alg:direct} (line 2) equals the matrix $G$ of \eqref{eq:system};
the right-hand side $b=\vecop(KB)=(I_r\otimes K)\vecop(B)$ by the identity
$\vecop(KB)=(I_r\otimes K)\vecop(B)$, and $B=TZ$ because
$B[i,:]=\sum_{\ell:i_\ell=i}t_\ell\tilde Z[\ell,:]=\sum_m T[i,m]Z[m,:]=(TZ)[i,:]$
(only observed $(i,m)$ have $T[i,m]=t_\ell\neq0$). By Proposition~\ref{prop:spd}
$G\succ0$, so Cholesky exists and the triangular solves return the unique
solution of $G\,w=b$.

\emph{(b)} The function \textsc{Matvec}$(W)$ returns $\mathcal G(W)$, the matrix
form of $G\vecop(W)$, by Lemma~\ref{lem:matvec}; \textsc{Prec}$(R)$ returns the
matrix form of $\mathcal M^{-1}\vecop(R)$ since $K=LL^{\mathsf T}$ implies the two
triangular solves yield $K^{-1}R$ columnwise, and $\mathcal M^{-1}=I_r\otimes
K^{-1}$ acts columnwise as $K^{-1}$. The scalars and updates in lines 13--22 are
exactly the preconditioned CG recurrences written in the matrix variables (with
$\langle\cdot,\cdot\rangle$ the Euclidean inner product of the vectorizations),
so the iterates coincide with PCG applied to $G\,\vecop(W)=b$ with SPD
preconditioner $\mathcal M$. Since $G\succ0$ and $\mathcal M\succ0$, PCG
converges; in exact arithmetic it computes the exact solution in at most
$\dim=nr$ steps because the $\mathcal M^{-1}G$-Krylov space stabilizes after at
most $nr$ iterations. The break condition makes the returned $W$ satisfy
$\langle R,R\rangle\le\tau^2\rho_0$, i.e. $\|b-G\vecop(W)\|_2\le\tau\|b\|_2$.
Letting $\tau\to0$ (or running $nr$ steps) recovers the exact solution of part
(a).
\end{proof}

\begin{proposition}[Termination]\label{prop:term}
Algorithm~\ref{alg:pcg} halts after at most $T_{\max}$ iterations; in exact
arithmetic with $T_{\max}\ge nr$ and $\tau=0$ it halts with the exact solution.
\end{proposition}
\begin{proof}
The \texttt{for} loop runs at most $T_{\max}$ times by construction. The exact
arithmetic claim is the finite termination of CG used in
Theorem~\ref{thm:correct}(b).
\end{proof}

\section{Complexity}

\begin{theorem}\label{thm:complexity}
Let $\kappa=\mathrm{cond}_2(\mathcal M^{-1}G)$ and let $\mathcal I$ be the number
of PCG iterations actually performed.
\begin{enumerate}[label=(\alph*)]
\item One PCG iteration costs $O(n^2r+qr)$ time (Lemma~\ref{lem:matvec} matvec
$O(n^2r+qr)$, preconditioner solve $O(n^2r)$, inner products/axpys $O(nr)$).
\item Total time of Algorithm~\ref{alg:pcg} is
\[
  O\!\big(\,\underbrace{n^3}_{\text{Cholesky}}+\underbrace{qr+n^2r}_{\text{rhs}}
  +\mathcal I\,(n^2r+qr)\,\big),
\]
with memory $O(nr+q(n+r))$. Since CG attains residual reduction $\epsilon$ in
$\mathcal I=O(\sqrt{\kappa}\,\log(1/\epsilon))$ iterations, the time is
$O\!\big(n^3+ \sqrt{\kappa}\,\log(1/\epsilon)\,(n^2r+qr)\big)$.
\item Algorithm~\ref{alg:direct} costs $O(q\,n^2r^2+n^3r^3)$ time and
$O(n^2r^2)$ memory.
\end{enumerate}
In particular, whenever $\mathcal I\,(n^2r+qr)\ll n^3r^3$ — which holds because
$\mathcal I$ is small once $\mathcal M=I_r\otimes K$ removes the kernel
conditioning, and $q,n,r$ are all $\ll$ the size of $G$ raised to its powers —
PCG is asymptotically cheaper, and no step touches a quantity of order $M$ or
$N$.
\end{theorem}

\begin{proof}
\emph{(a)} Each line of \textsc{Matvec} was costed after Lemma~\ref{lem:matvec}:
$Y=KW$ is $O(n^2r)$, computing all $u_\ell$ is $O(qr)$, the scatter into $C$ is
$O(qr)$, and the final $K(C+\lambda W)$ is $O(n^2r)$; total $O(n^2r+qr)$.
\textsc{Prec} is $2r$ triangular solves with $n\times n$ factors, $O(n^2r)$. The
vector operations on $n\times r$ matrices are $O(nr)$. Summing gives $O(n^2r+qr)$.

\emph{(b)} The precomputation is one $n\times n$ Cholesky $O(n^3)$ plus the RHS
$O(qr+n^2r)$. The loop runs $\mathcal I$ times at $O(n^2r+qr)$ each, by (a).
Memory holds $\tilde Z$ ($O(qr)$), the indices and $\tilde K$/values ($O(qn)$ or
$O(q)$), $K$ and $L$ ($O(n^2)$), and a constant number of $n\times r$ work
matrices ($O(nr)$): total $O(nr+q(n+r)+n^2)=O(nr+q(n+r))$ since $n\le q$. The
classical CG bound
$\|x_{\mathcal I}-x_\star\|_{\mathcal M^{-1}G}\le
2\big(\tfrac{\sqrt\kappa-1}{\sqrt\kappa+1}\big)^{\mathcal I}\|x_0-x_\star\|_{\mathcal M^{-1}G}$
gives $\mathcal I=O(\sqrt\kappa\log(1/\epsilon))$.

\emph{(c)} Line 1 of Algorithm~\ref{alg:direct} builds $q$ rows of length $nr$:
$O(q\,nr)$. Line 2 forms the $nr\times nr$ Gram $P^{\mathsf T}P$ in
$O(q(nr)^2)=O(q\,n^2r^2)$ and adds the regularizer in $O(n^2r)$. Line 4 is a
dense Cholesky of an $nr\times nr$ matrix, $O((nr)^3)=O(n^3r^3)$. Lines 3,5,6 are
lower order. Memory is dominated by storing $G$, $O(n^2r^2)$.
\end{proof}

\paragraph{Discussion of the speedup.}
The direct cost $O(q\,n^2r^2+n^3r^3)$ has the cubic-in-$nr$ factor the statement
flags. PCG replaces it by $O(n^3)+\mathcal I\cdot O(n^2r+qr)$. The one-off
$O(n^3)$ Cholesky of the \emph{small} $n\times n$ kernel ($n\ll M$) is negligible,
and each iteration is linear in $q$ and in $r$ and quadratic only in the small
$n$. The preconditioner $\mathcal M=I_r\otimes K$ is precisely what keeps
$\mathcal I$ small: it maps the regularizer block exactly to $\lambda I$ and
strips the (typically catastrophic) kernel ill-conditioning out of the data
block, leaving a spectrum clustered around $\lambda$ plus a low-rank
($\le q$) data perturbation. All of this avoids forming $Z\otimes K$
($M n\times r n$), $S$ ($N\times q$), or any length-$M$/length-$N$ vector; the
only objects of size beyond $O(n^2)$ are $\tilde Z\in\bR^{q\times r}$ and the
$q$ index/value triples, all of size $O(q(n+r))\ll N$.

\section*{Numerical verification (summary)}

The matvec identities \eqref{eq:Pw}--\eqref{eq:matvec} and the factorization
\eqref{eq:G-PtP} were checked against a dense assembly of $G$ on random
instances (agreement to machine precision). On a smooth squared-exponential
kernel with $n=20$, $r=4$, $q=300$, $\lambda=10^{-2}$, the kernel had
$\mathrm{cond}(K)\approx10^{11}$ and $\mathrm{cond}(G)\approx2\times10^{15}$;
the preconditioner $\mathcal M=I_r\otimes K$ reduced
$\mathrm{cond}(\mathcal M^{-1}G)$ to $\approx2\times10^{4}$. Matrix-free PCG
matched the dense direct solve and converged in $96$ iterations versus $563$ for
unpreconditioned CG (about a $6\times$ reduction), confirming both correctness
and the effectiveness of the preconditioner.

\end{document}
